\section{Trazabilidad de requisitos}
\label{sec:trazabilidad}
\textit{Responsable: Mariano.}

\begin{longtable}{p{0.40\textwidth}p{0.12\textwidth}p{0.38\textwidth}}
\toprule
Requisito del enunciado & Sección & Cómo se cumple \\
\midrule
\endhead
SaaS multi-tenant & \ref{sec:seguridad} & \texttt{tenant\_id} en el token, en cada evento y en la clave de cada tabla. \\
Serverless & \ref{sec:arquitectura} & API Gateway + Lambda, DynamoDB, EventBridge, SQS y SNS. \\
Basada en eventos & \ref{sec:eventos} & Outbox + Pipes $\rightarrow$ EventBridge $\rightarrow$ una SQS por consumidor; SNS FIFO para notas. \\
Mínimo 5 microservicios protegidos & \ref{sec:microservicios}, \ref{sec:seguridad} & 6 microservicios, cada uno con su API Gateway y authorizer de Cognito. \\
Base de datos propia por microservicio & \ref{sec:microservicios} & Una tabla DynamoDB por microservicio (ADR-007). \\
Estructura JSON de cada base & \ref{sec:microservicios} & Subsección de cada microservicio. \\
30\,000 registros en 3 tenants & \ref{sec:microservicios} & \texttt{poblar\_datos.py}: UTEC 15\,000, UPC 10\,000, PUCP 5\,000. \\
Uso de S3, SNS y SQS & \ref{sec:microservicios}, \ref{sec:eventos} & S3 para archivos; SNS FIFO y correo; SQS por consumidor. \\
Despliegue con Serverless Framework & \ref{sec:despliegue} & Un \texttt{serverless.yml} por stack y \texttt{desplegar.sh}. \\
Frontend en Amplify & \ref{sec:frontend} & \texttt{infra/frontend} y \texttt{nubeduca-frontend}. \\
Diagrama de arquitectura & \ref{sec:arquitectura} & Figura~\ref{fig:arquitectura}. \\
Demo con 2 tenants & \ref{sec:pruebas}, \ref{sec:repos} & QA-01 con UTEC y UPC; video en \url{https://drive.google.com/drive/folders/1vRckiyNb-z5vK4JJRX3VzP7CAGKojie-}. \\
Enlaces a los repositorios & \ref{sec:repos} & Tabla de repositorios públicos. \\
\bottomrule
\end{longtable}
